\subsection{A bounded character-sum attempt at the short unit window}

This is a check of a proposed route to the missing simultaneous
unit, not a new existence theorem.  Fix $p\geq31$ prime with
$p\equiv3\pmod4$, put $b=(p+3)/2$, and let $D=-d$ be odd,
negative, fundamental and prime to $p$.  Write
$\chi=\chi_D$, $\epsilon(r)=(\frac rp)$, and
$P(j)=\sum_{0\leq a<j}\chi(a)$ for $0\leq j<d$.
The sufficient simultaneous-unit window already proved is
$d^5\leq1600p^2$.  Within it the third value $\beta_3(\chi)$
and the class number are units; the missing assertion is a
unit of $\beta_b(\chi)$ for at least one such $D$.

\paragraph{The exact integer representative.}
KT-FIELD-CHARACTER-REDUCTION and the complete prefix formula give
\[
 \beta_b(\chi)\equiv I_D
 :=\chi(p)\sum_{r=1}^{p-1}r\epsilon(r)
                         P([p^{-1}r]_d)\pmod p.
 \tag{WC.1}
\]
Indeed $b-1=(p-1)/2+1$, so Euler's criterion changes
$r^{b-1}$ into $r\epsilon(r)$ in that formula.
This is a signed integer representative of the residue;
it is not an equality of ordinary real values with
$\beta_b(\chi)$.

For clarity, assume $d<p$, as holds in the proposed short
window.  For $0\leq j<d$ let $a_j\in\{1,\ldots,d\}$
represent $pj$ modulo $d$, set
\[
 N_j=1+\left\lfloor\frac{p-1-a_j}{d}\right\rfloor,
 \qquad u_j\equiv a_jd^{-1}\pmod p,
 \quad 0\leq u_j<p,
\]
and put
\[
 S_j=\sum_{t=0}^{N_j-1}\epsilon(u_j+t),\qquad
 T_j=\sum_{t=0}^{N_j-1}t\epsilon(u_j+t).
\]
Partitioning $1\leq r<p$ according to $r\equiv pj\pmod d$
gives the exact identity of integers
\[
 I_D=\chi(p)\epsilon(d)
        \sum_{j=0}^{d-1}P(j)(a_jS_j+dT_j).
 \tag{WC.2}
\]
The shifts $u_j$ can be anywhere modulo $p$; the character
is periodic, so the intervals may pass through $p$.
Their lengths are about $p/d$.  Formula (WC.2) has
$d$ intervals, their prefix coefficients, and the factors
$a_j,d$.  A bound for one unweighted character sum is
not itself a bound below $p$ for $I_D$.

\paragraph{What the direct Burgess estimate actually gives.}
For a shifted interval of length $X$ and weights
$t+\alpha$, $0<\alpha\leq1$, summation by parts and
the classical Burgess estimate give, for every integer
$R\geq1$,
\[
 \left|\sum_{0\leq t<X}(t+\alpha)\epsilon(u+t)\right|
 \ll_{R,\varepsilon}
       X^{2-1/R}p^{(R+1)/(4R^2)+\varepsilon}.
 \tag{WC.3}
\]
Primality of $p$ satisfies the cube-free modulus condition
for arbitrary $R$.  The normalized expression for a block
in (WC.2) is of this form with $\alpha=a_j/d$; recovering
that block multiplies by $d$.

Even discard this last factor, all prefix coefficients,
and the sum over $j$, which is more favorable than the
actual expression.  At the edge $d\asymp p^{2/5}$ of
the desired window one has $X\asymp p^{3/5}$, and the
power of $p$ in this one-block envelope is
\[
 E_R=\frac65-\frac7{20R}+\frac1{4R^2}.
 \tag{WC.4}
\]
Its minimum over positive integer $R$ is
$E_2=87/80>1$.  In fact $E_1=11/10$, and $E_R$
is increasing for real $R\geq2$ since
$E_R'=(7R-10)/(20R^3)>0$ there.
Thus this estimate does not bound even the favorable
normalized block by $p$ in the proposed range.
The constants and logarithmic factors cannot change
that asymptotic power obstruction.

More generally, for $d\asymp p^\alpha$ the same envelope
has exponent below $1$ only when
\[
 \alpha>
 \frac{4R^2-3R+1}{4R(2R-1)}.
 \tag{WC.5}
\]
The smallest right side over positive integer $R$ is
$11/24$, at $R=2$.  For $R\geq2$ its difference from
$11/24$ is
$(2R-3)(R-2)/(24R(2R-1))\geq0$, while $R=1$
gives $1/2$.  The short-unit exponent $2/5$ is smaller.
These calculations delimit this direct absolute-value
argument; they are not lower bounds on the actual sums
and do not exclude additional cancellation between blocks.

\paragraph{Why the positive fixed-weight value is not a small substitute.}
The already checked fixed-weight identity is
\[
 \beta_b(\chi_D)\equiv\beta_2(\chi_F)\pmod p,
 \qquad F=\mathbf Q(\sqrt{pd}),\quad
 \chi_F=\chi_D\chi_{-p}.
 \tag{WC.6}
\]
Put $f=pd$.  The elementary Bernoulli Fourier expansion,
with the positive quadratic Gauss sum $\sqrt f$, gives
\[
 \beta_2(\chi_F)=\frac{f^{3/2}}{2\pi^2}L(2,\chi_F)>0.
 \tag{WC.7}
\]
Every Euler factor of this positive $L$-value is at least
$(1+\ell^{-2})^{-1}$, including primes dividing $f$.
Consequently
\[
 L(2,\chi_F)\geq\frac{\zeta(4)}{\zeta(2)}
                  =\frac{\pi^2}{15},\qquad
 \beta_2(\chi_F)\geq\frac{(pd)^{3/2}}{30}>p
       \quad(p\geq43,\ d\geq3).
 \tag{WC.8}
\]
The final inequality follows from $27p>900$.
So positivity of this real value cannot prove that it
is a nonzero number smaller than $p$.  Passing instead
to the smaller residue representative (WC.1) loses
that positivity and requires its own nonvanishing argument.

The existence of a small quadratic nonresidue likewise
does not by itself supply a nonzero weighted moment.
For example, modulo $31$ the first three Legendre signs
are $1,1,-1$, but $1+2-3=0$.  More globally the
ordinary class-number formula gives
$\sum_{r=1}^{p-1}r\epsilon(r)=-p\,h(-p)$ for
$p>3$, $p\equiv3\pmod4$: a natural nonzero weighted
sum can be divisible by $p$.  These examples concern
the proposed inference from character signs, not a
counterexample to the target simultaneous field existence.

\paragraph{Precise remaining requirement and source scope.}
One would need a nonzero residue in (WC.1) for some
fundamental $d$ in the short window, with a method that
retains the same $D$ for all units.  Separate real
nonvanishing, a bound for the least nonresidue, or an
upper bound for one block in (WC.2) does not prove this.
No such additional cancellation or lower-bound statement
is established in this attempt.

The analytic input is the classical Burgess statement
as stated in E.~Trevi\~no, \emph{The Burgess inequality
and the least $k$-th power non-residue}, primary author
PDF \url{https://campus.lakeforest.edu/trevino/Burgess.pdf},
pp.~1--2.  The full statement, including the cube-free
qualification, and his explicit Theorem~1 and Corollary~1
on p.~2 were read.  Only these estimate scopes are
imported; the foundational Burgess proof was not newly
reviewed.  The Fourier calculation and exponent
optimization above are displayed deductions.
There is no broad numerical scan, no new main-gate
theorem, and no claim that every character-sum method
must fail.
